\documentclass[10pt]{article}
\usepackage[T1]{fontenc}
\usepackage[utf8]{inputenc}
\usepackage[margin=1in]{geometry}
\usepackage{amsmath,amssymb,amsthm,booktabs,array}
\usepackage[colorlinks=true,urlcolor=blue,citecolor=blue,linkcolor=blue]{hyperref}
\newtheorem{theorem}{Theorem}
\newtheorem{lemma}[theorem]{Lemma}
\newtheorem{proposition}[theorem]{Proposition}
\newtheorem{corollary}[theorem]{Corollary}
\newcommand{\FF}{\mathbb F_2}
\newcommand{\WW}{\mathcal W}
\newcommand{\LL}{\mathcal L}
\title{A degree-12 counterexample to Saniga et al.'s\newline Section 4.4 line-subgeometry conjecture}
\author{Haobo Ma\\{\small ChronoAI Pte Ltd, Singapore}
\and Wenlin Zhang\\{\small National University of Singapore}\\{\small ChronoAI Pte Ltd, Singapore}}
\date{2 November 2026}
\begin{document}
\maketitle
\begin{abstract}
We refute the line-subgeometry conjecture in Section~4.4, p.~20 of Saniga et al.,
\emph{Hexagons govern three-qubit contextuality}, arXiv:2312.07738v3
(\emph{Quantum} \textbf{9} (2025), 1601).
The configuration of all 75 isotropic lines contained in a point's perp
hyperplane has contextuality degree exactly 12, whereas it shares exactly 15
lines with every classically embedded split Cayley hexagon.
The intersection count requires only three lines through each ambient point.
The degree follows from a general identity: for every nonzero pole in $n\ge2$
qubits, the full contained-line geometry of its perp has four times the degree
of the full $(n-1)$-qubit line geometry; in four qubits this gives degree 252.
We also give a 19-line counterexample of degree one.
\end{abstract}

\section{The conjecture and the configuration}
Saniga et al.~\cite[Section~4.4, p.~20]{saniga} propose a generalization of their
hexagon description of contextuality for the whole three-qubit space, quadrics,
and doilies: ``A natural generalization of these findings is that this should
hold for any contextual sub-geometry of $\WW(5,2)$ whose contexts are lines.''
In particular, their conjecture would imply that every contextual line family
$C$ admits a classically embedded split Cayley hexagon $H$ with
\begin{equation}\label{eq:conjecture}
 d(C)=|C\cap H|.
\end{equation}
Here a hexagon is identified with its line set. We disprove even this numerical
consequence; consequently the stronger proposed description of optimal violated
contexts also fails in this generality.
The established special-case results are unaffected.

Let $V_n=(\FF^2)^n$. Write a label as
$v=((x_1,z_1),\ldots,(x_n,z_n))$ and put
\[
 B(u,v)=\sum_{j=1}^n(x_jz'_j+z_jx'_j),\qquad
 P_v=\bigotimes_{j=1}^n i^{x_jz_j}X^{x_j}Z^{z_j}.
\]
Thus the local labels $(0,0),(1,0),(0,1),(1,1)$ denote $I,X,Z,Y$,
respectively. The points of $\WW(2n-1,2)$ are the nonzero labels; its
isotropic lines are $\ell=\{u,v,u+v\}$ with $u,v$ distinct and nonzero and
$B(u,v)=0$. Denote their set by $\LL_n$. Define $s(\ell)\in\FF$ by
$\prod_{q\in\ell}P_q=(-1)^{s(\ell)}I$ and set
\[
 d(C)=\min_{f:V_n\setminus\{0\}\to\FF}
       \bigl|\{\ell\in C:\textstyle\sum_{q\in\ell}f(q)\ne s(\ell)\}\bigr|.
\]
The product is order-independent since the three observables commute.
All minima exist because the assignment sets are finite.

For a nonzero $p\in V_3$, let
\[
 S_p=\{q\ne0:B(p,q)=0\},\qquad
 C_p=\{\ell\in\LL_3:\ell\subseteq S_p\}.
\]
The family $C_p$ includes lines not through $p$. It is therefore different
from the perpset in Muller et al.~\cite[Section~2 and Proposition~2]{muller},
whose contexts are only the lines through the pole and whose degree is zero.
The finite obstruction $0<d(C_{IIZ})\le12<15\le|C_{IIZ}\cap H|$ is also
encoded in the pinned Lean declaration described in Section~\ref{sec:scope}.
The arguments below determine both numbers exactly and transport the degree
to every pole.

\begin{theorem}\label{thm:main}
For every nonzero $p\in V_3$, the configuration $C_p$ has 31 points, 75 lines,
and contextuality degree 12. If $H\subseteq\LL_3$ has exactly three lines
through each of the 63 ambient points, then $|C_p\cap H|=15$.
In particular, \eqref{eq:conjecture} fails for $C_p$ and every classically
embedded split Cayley hexagon.
\end{theorem}

\section{An exact degree by linear sections}
Put $d_m=d(\LL_m)$. The following identity supplies the lower bound without
enumerating assignments on 31 points.
\begin{proposition}\label{prop:lift}
For $n\ge2$ and every nonzero $p\in V_n$, let $C^{(n)}_p$ contain all isotropic
lines in $p^\perp\setminus\{0\}$. Then $d(C^{(n)}_p)=4d_{n-1}$.
\end{proposition}
\begin{proof}
First take $p=I^{\otimes(n-1)}Z$ and write $m=n-1$.
Every label in $p^\perp$ has the form $(u,t)$ with
$u\in V_m$ and $t\in\FF$, representing $P_u\otimes Z^t$; the pole is
$(0,1)$. There are two types of lines. The lines through the pole are
\[
 \{(0,1),(u,0),(u,1)\},\qquad u\in V_m\setminus\{0\}.
\]
Their sign is positive. Every other line projects injectively to a line
$\{u,v,u+v\}\in\LL_m$ and has the form
\[
 \{(u,a),(v,b),(u+v,a+b)\},\qquad a,b\in\FF.
\]
These are four distinct lifts, all with the sign of the base line because
$Z^aZ^bZ^{a+b}=I$. This also proves that the displayed list is exhaustive:
a line whose projection is not injective contains the pole, and otherwise
its projected labels are distinct, nonzero, and sum to zero.

For each linear functional $\lambda:V_m\to\FF$, the section
$u\mapsto(u,\lambda(u))$ identifies a copy of the full signed geometry
$\LL_m$ inside $C^{(n)}_p$. Any assignment on the perp violates at least
$d_m$ lines in each section. There are $2^{2m}$ sections. A fixed non-pole
line prescribes $\lambda(u)=a$ and $\lambda(v)=b$ on two independent
vectors, and so belongs to exactly $2^{2m-2}$ sections. Pole lines belong
to none. If $N$ is the number of violated non-pole lines, double counting gives
\[
 2^{2m-2}N\ge2^{2m}d_m,\qquad N\ge4d_m.
\]
This remains valid for $m=1$, when the base line set is empty and $d_1=0$.

Conversely, choose an assignment $g$ attaining $d_m$ and put
$f(u,t)=g(u)$ for $u\ne0$ and $f(0,1)=0$. Every pole line is satisfied,
and each base violation gives exactly four lifted violations. Thus
$d(C^{(n)}_p)\le4d_m$, proving equality.

For an arbitrary nonzero pole $p$, single-qubit Clifford gates first change
every nonidentity factor of $P_p$ to $Z$; controlled-NOT gates reduce a product
of $Z$ factors to one, and swaps move it to the last qubit. Thus a Clifford
unitary induces a symplectic map $g$ taking $p$ to $I^{\otimes(n-1)}Z$.
Write $UP_vU^\dagger=(-1)^{\tau(v)}P_{gv}$. Conjugating the product on a line
gives
\[
 s(g\ell)=s(\ell)+\sum_{v\in\ell}\tau(v).
\]
The assignment correspondence $f'(gv)=f(v)+\tau(v)$ therefore preserves the
violation status of every line. Since $g$ carries $C^{(n)}_p$ bijectively
onto the coordinate perp geometry, their degrees agree. This is the signed
Clifford transport principle of~\cite[Lemma~3]{li}, with the assignment
correspondence made explicit here.
\end{proof}

\begin{lemma}\label{lem:doily}
The full two-qubit line geometry has degree $d_2=3$.
\end{lemma}
\begin{proof}
For completeness, Tables~\ref{tab:lines} and~\ref{tab:grids} give an explicit
parity certificate for the familiar doily value
\cite[Section~2, Figure~4]{saniga}. A line has sign $-1$ precisely for
numbers 10, 13, and 14 in Table~\ref{tab:lines}; for example
$(XX)(ZZ)(YY)=-I$. The table follows directly from
$XY=iZ$, $YZ=iX$, $ZX=iY$, reversing the sign when two distinct
nonidentity factors are interchanged.
Each row of Table~\ref{tab:grids} consists of six lines in which every
occurring point appears exactly twice and the number of negative lines is odd.
Summing the six parity equations makes their left side zero and their right
side one, so at least one line in each row is violated. Each of the 15 lines
occurs in four rows. If an assignment violates $k$ lines, then $4k\ge10$,
and integrality gives $k\ge3$. The zero assignment violates exactly the
three negative lines, proving $d_2=3$.
This ten-grid incidence argument is also the local lower-bound method used
by Li, Gupta, and Ganesh~\cite[Lemma~8 and Proposition~10]{li}.
\end{proof}

\begin{table}[ht]
\centering\small\setlength{\tabcolsep}{3pt}
\begin{tabular}{clc@{\qquad}clc@{\qquad}clc}
\toprule
No.&Observables&Sign&No.&Observables&Sign&No.&Observables&Sign\\
\midrule
1&$XI,IX,XX$&$+$&6&$ZI,IY,ZY$&$+$&11&$XX,YZ,ZY$&$+$\\
2&$XI,IZ,XZ$&$+$&7&$YI,IX,YX$&$+$&12&$ZX,XZ,YY$&$+$\\
3&$XI,IY,XY$&$+$&8&$YI,IZ,YZ$&$+$&13&$ZX,YZ,XY$&$-$\\
4&$ZI,IX,ZX$&$+$&9&$YI,IY,YY$&$+$&14&$YX,XZ,ZY$&$-$\\
5&$ZI,IZ,ZZ$&$+$&10&$XX,ZZ,YY$&$-$&15&$YX,ZZ,XY$&$+$\\
\bottomrule
\end{tabular}
\caption{All two-qubit isotropic lines, with their canonical Pauli signs.}
\label{tab:lines}
\end{table}
\begin{table}[ht]
\centering\small
\begin{tabular}{c@{\quad}c}
\toprule
Six line numbers&Six line numbers\\
\midrule
$1,2,4,5,10,12$&$2,3,8,9,12,13$\\
$1,2,7,8,11,14$&$4,5,7,8,13,15$\\
$1,3,4,6,11,13$&$4,6,7,9,12,14$\\
$1,3,7,9,10,15$&$5,6,8,9,10,11$\\
$2,3,5,6,14,15$&$10,11,12,13,14,15$\\
\bottomrule
\end{tabular}
\caption{Ten parity-contradictory grids; every line occurs four times.}
\label{tab:grids}
\end{table}

\begin{proof}[Proof of the degree and size assertions in Theorem~\ref{thm:main}]
Proposition~\ref{prop:lift} and Lemma~\ref{lem:doily} give
$d(C_p)=12$. At $IIZ$ the 31 points are the pole and two lifts of each base point.
The 15 pole lines and 60 other lines give 75 lines in total;
the same counts hold at every pole by the bijection in that proposition.
The degree is pole-independent, but the number of negative lines for the
canonical Pauli representatives is not: direct multiplication gives 12 at
$IIZ$ and 24 at $YYY$, as the accompanying finite calculation verifies.
At $IIZ$ the zero assignment violates exactly the four lifts of each of
lines 10, 13, and 14.
\end{proof}

\begin{corollary}\label{cor:optimal}
An assignment on $C_{IIZ}$ is optimal if and only if all 15 pole lines are
satisfied and each of the 16 linear sections has exactly three violated lines.
\end{corollary}
\begin{proof}
Every section has at least three violations. If there are 12 violations in
total, the section count is at least $16\cdot3=48$ and at most
$4\cdot12=48$. Equality forces every section to have three and every
violation to avoid the pole. Conversely, the stated conditions give 48
section incidences and hence $48/4=12$ total violations.
\end{proof}

\section{The intersection obstruction and its scope}
\begin{proof}[Proof of the intersection assertion in Theorem~\ref{thm:main}]
The nonzero linear functional $B(p,-)$ has a 32-element kernel, so
$|S_p|=31$. On any isotropic line, its restriction is either zero everywhere
or has a one-dimensional kernel. Thus every line meets $S_p$ in three or
one points, respectively. If $H$ has degree three at every ambient point,
counting all incidences gives $3|H|=3\cdot63$, so $|H|=63$.
Put $k=|C_p\cap H|$. Counting incidences of $H$ with $S_p$ gives
\[
 3\cdot31=3k+(63-k),\qquad k=15.
\]
A classically embedded split Cayley hexagon has all 63 ambient points and
three isotropic lines through each point
\cite[Section~2, pp.~6--7]{saniga}. It therefore satisfies the hypotheses,
and $d(C_p)=12<15=|C_p\cap H|$ proves the contradiction.
\end{proof}

\paragraph{A surviving upper bound.}
For any signed line family $C$ and a hexagon $H$ whose full line complement
admits a satisfying assignment, that assignment can violate only lines of
$C\cap H$. Hence $d(C)\le|C\cap H|$.
Complements of classically embedded hexagons are noncontextual
\cite{holweck}; this supplies the premise for genuine classical embeddings.
Theorem~\ref{thm:main} shows that this upper bound can be strict for every
such embedding simultaneously. Additional hypotheses are necessary for
equality.

\begin{corollary}\label{cor:small}
There is a 19-line family $C$ with $d(C)=1$ and $|C\cap H|\ge3$ for every
three-regular isotropic line family $H$ on all 63 points.
\end{corollary}
\begin{proof}
Take all 15 isotropic lines through $XII$ and add the rows and columns of
the Mermin--Peres square
\[
\begin{array}{ccc}
XII&IXI&XXI\\ IZI&ZII&ZZI\\ XZI&ZXI&YYI.
\end{array}
\]
Exactly two square lines contain $XII$, so the union has $15+6-2=19$
lines. All lines through $XII$ are positive: their other two observables
differ only by an $X$ on the first qubit and their product is $XII$.
The square has just one negative line, $\{XXI,ZZI,YYI\}$.
Thus the zero assignment has one violation, while the square's odd parity
forces at least one. Finally, $H$ contains three lines through $XII$, all in
$C$, so $|C\cap H|\ge3$.
\end{proof}
\paragraph{Consequences and neighboring statements.}
Proposition~\ref{prop:lift} gives the exact degree $4d_{n-1}$ at every nonzero
pole in $n\ge2$ qubits, not just at a coordinate pole. In particular, it gives
degree $252$ for every full contained-line perp geometry in four qubits,
using the known value
$d_3=63$~\cite[Section~4.1]{saniga}. This is a statement about that perp,
not the degree of the full four-qubit space. Table~\ref{tab:coverage}
separates the conjecture from the neighboring assertions in its source.

\begin{table}[ht]
\centering\small
\begin{tabular}{p{.43\linewidth}p{.50\linewidth}}
\toprule
Source item in arXiv:2312.07738v3&Mathematical coverage\\
\midrule
Section~4.4, p.~20: arbitrary contextual line subgeometries&
Refuted by $C_p$: degree 12, intersection 15.\\
Sections~4.1 and 4.4: full space, elliptic and hyperbolic quadrics, doilies&
Restricted assertions retained; the coordinate doily degree 3 is reproved above.\\
Section~4.4, pp.~20--21: skew-hexagon complements (24-line result on p.~21)&
Not settled by the perp argument; the source's restricted result is distinct.\\
Section~6, p.~24: four-qubit hyperbolic quadric&
Degree 315 and classification of the 240 optimal violated sets as copies of
$DW(5,2)$ proved subsequently by Li et al., Theorems~11 and~15~\cite{li};
the classification uses a DRAT-certified finite enumeration.\\
Section~6, pp.~24--25: full four-qubit space and proposed higher analogue&
Not settled here; neither the perp formula nor the quadric theorem supplies the proposed global configuration.\\
\bottomrule
\end{tabular}
\caption{Coverage of the conjecture and its neighboring source statements.}
\label{tab:coverage}
\end{table}

\section{Formal statements and finite certificates}\label{sec:scope}
The immutable \href{https://github.com/the-omega-institute/trureturing/blob/6715a5bea38a03c80f6c4da6de93fea88798f026/D5/S3/Quantum/Contextuality/PerpHyperplaneHexagonContextualityRefutation.lean}{Lean source}
is
\begin{quote}\small\ttfamily
D5/S3/Quantum/Contextuality/\\
PerpHyperplaneHexagonContextualityRefutation.lean
\end{quote}
It defines \texttt{claim} and proves \texttt{result : $\neg$ claim}, with qualified names
\begin{quote}\small\ttfamily
D5.S3.Quantum.Contextuality.\\
PerpHyperplaneHexagonContextualityRefutation.claim\\[2pt]
D5.S3.Quantum.Contextuality.\\
PerpHyperplaneHexagonContextualityRefutation.result
\end{quote}
The hyperlink fixes the source commit; \texttt{verification-scope.json} gives its full
identifier, both qualified declaration names, and the version and digest of
\texttt{verify.py}.
Its \texttt{IsClassicalHexagon} predicate requires isotropic lines, degree three
at every ambient point, no ordinary triangles, and planarity at every point.
Every genuine classical embedding satisfies these conditions by the
definitions in~\cite[Sections~2--3]{saniga}; the incidence proof above only
needs isotropy and degree three at every ambient point.

The declaration's proof uses the fixed pole $IIZ$, positivity of its degree,
the upper bound 12, and the lower bound 15 on hexagon intersections.
The exact degree, the exact intersection count for arbitrary three-regular
line families, the transport to all poles, Proposition~\ref{prop:lift}, and
Corollary~\ref{cor:optimal} have the written proofs given here.
Corollary~\ref{cor:small} is likewise a written proof.
The accompanying \texttt{verify.py} checks all 315 three-qubit line signs by
two implementations (matrix multiplication and binary phase exponents), the
15-line and ten-grid tables, every one of the $2^{15}$ two-qubit assignments,
the 16 sections and their multiplicities, and the perp's 31 points and 75 lines.
It also checks the intersection count on the classical coordinate model from
\cite[equations~(13)--(14)]{saniga}.
The pole-dependent negative-line counts are checked for all 63 poles.
The arbitrary-$n$, arbitrary-pole identity is proved by
the section argument, rather than inferred from these finite checks.
The 19-line union, its signs, and its square parity are checked as well.

\begin{thebibliography}{9}
\bibitem{saniga}
M. Saniga, F. Holweck, C. Kelleher, A. Muller, A. Giorgetti, and H. de Boutray,
\emph{Hexagons govern three-qubit contextuality}, Quantum \textbf{9} (2025), 1601.
\href{https://doi.org/10.22331/q-2025-01-20-1601}{doi:10.22331/q-2025-01-20-1601};
\href{https://arxiv.org/abs/2312.07738v3}{arXiv:2312.07738v3}.
\bibitem{muller}
A. Muller, M. Saniga, A. Giorgetti, H. de Boutray, and F. Holweck,
\emph{New and improved bounds on the contextuality degree of multi-qubit configurations},
Mathematical Structures in Computer Science \textbf{34} (2024), 322--343.
\href{https://doi.org/10.1017/S0960129524000057}{doi:10.1017/S0960129524000057};
\href{https://arxiv.org/abs/2305.10225}{arXiv:2305.10225}.
\bibitem{holweck}
F. Holweck, H. de Boutray, and M. Saniga,
\emph{Three-qubit-embedded split Cayley hexagon is contextuality sensitive},
Scientific Reports \textbf{12} (2022), 8915.
\href{https://doi.org/10.1038/s41598-022-13079-3}{doi:10.1038/s41598-022-13079-3}.
\bibitem{li}
Z. Li, A. Gupta, and V. Ganesh,
\emph{The Four-Qubit Hyperbolic Quadric Has Contextuality Degree 315},
\href{https://arxiv.org/abs/2610.26465v1}{arXiv:2610.26465v1} (2026).
\end{thebibliography}
\enlargethispage{2\baselineskip}
\medskip
\noindent\small Licensed under Creative Commons Attribution 4.0 International
(\href{https://creativecommons.org/licenses/by/4.0/}{CC BY 4.0}).
\end{document}
